\uuid{bzPe}
\titre{Distance}
\auteur{bodin ; hanani ; mzari}
\organisation{exo7}
\licence{CC BY-NC-SA 4.0}
\datecreate{2020-01-20}
\niveau{L1}
\module{Géométrie}
\chapitre{Géométrie de l'espace}
\sousChapitre{Distance}
\difficulte{2}
\theme{Distance}
\exo7qcmid{lille-1.geometrie-espace.025}

\contenu{
\texte{
Pour ces questions, l'espace est muni d'un repère orthonormé direct  $(O, \vec{i}, \vec{j}, \vec{k})$.
}

\begin{question}
\qclassification{141.05}
\qtags{motcle=distance}
Soit $D$ la droite passant par le point $A(1,-1,0)$ et dirigée par le vecteur $\vec{u}(1,1,-1)$. Soit $M(1,-1,3)$ un point et $H$ le projeté orthogonal de $M$ sur $D$. Les coordonnées de $H$ sont :
\begin{answers}

    \bad {$H(0,1,1)$}

    \bad {$H(1,2,1)$}

   \good {$H(0,-2,1)$}

   \bad {$H(1,-2,1)$}

\end{answers}
\begin{explanations}
$H\in D$, donc il existe $t \in \Rr$ tel que $H(1+t,-1+t,-t)$. On calcule $t$ en utilisant l'égalité : $\overrightarrow{HM} \cdot \vec{u} = 0$.
\end{explanations}
\end{question}
}
